\documentclass[12pt]{article}
\usepackage{amsmath, amssymb, geometry, enumitem}
\geometry{margin=1in}
\setlength{\parindent}{0pt}
\setlength{\parskip}{0.5em}

\title{ECON 101: Macroeconomics \\ Quiz 5}
\author{Instructor: Muhebullah Karimzada}
\date{April 07,  2026}

\begin{document}
\maketitle



\begin{enumerate}[label=\textbf{\arabic*.}]

\item The primary monetary policy tool used by the Bank of Canada is:
\begin{enumerate}[label=(\Alph*)]
    \item The money supply target
    \item The exchange rate
    \item The target overnight interest rate
    \item Government spending
\end{enumerate}

\item A decrease in the target overnight rate will most directly:
\begin{enumerate}[label=(\Alph*)]
    \item Increase long-run aggregate supply
    \item Increase short-term interest rates
    \item Reduce borrowing costs and stimulate spending
    \item Decrease inflation immediately to zero
\end{enumerate}

\item Which of the following is part of the monetary transmission mechanism?
\begin{enumerate}[label=(\Alph*)]
    \item Changes in tax rates
    \item Changes in interest rates affecting investment
    \item Changes in labour laws
    \item Changes in population growth
\end{enumerate}

\item If the Bank of Canada raises interest rates, the Canadian dollar will most likely:
\begin{enumerate}[label=(\Alph*)]
    \item Depreciate
    \item Appreciate
    \item Remain unchanged
    \item Become fixed
\end{enumerate}
\newpage
\item Contractionary monetary policy will:
\begin{enumerate}[label=(\Alph*)]
    \item Shift AD to the right
    \item Increase investment spending
    \item Shift AD to the left
    \item Increase net exports
\end{enumerate}

\end{enumerate}



\textbf{Problem}

The Bank of Canada is concerned that inflation is above its target. It decides to increase the policy interest rate.

Suppose the economy is initially at the following equilibrium:
\begin{itemize}
    \item Potential GDP: \(Y^* = 2{,}500\)
    \item Actual GDP: \(Y = 2{,}700\)
    \item Inflation is rising
\end{itemize}

The Bank increases the policy interest rate, which reduces investment by \$100 billion.

The marginal propensity to consume (MPC) is \(0.75\).

\begin{enumerate}

\item[(a)] Calculate the spending multiplier. (3 marks)

\item[(b)] Calculate the total change in equilibrium GDP resulting from the decrease in investment. (4 marks)

\item[(c)] Calculate the new level of equilibrium GDP. (2 marks)

\item[(d)] Using an AD--AS diagram, illustrate:
\begin{itemize}
    \item the initial equilibrium,
    \item the shift in aggregate demand,
    \item the new short-run equilibrium.
\end{itemize}
Clearly label:
\begin{itemize}
    \item AD$_0$, AD$_1$,
    \item SRAS, LRAS,
    \item equilibrium output and price level.
\end{itemize}
(3 marks)

\item[(e)] Explain how this monetary policy helps reduce inflation in the medium run. (3 marks)

\end{enumerate}

\hrule
\vspace{1em}

\begin{center}
\textit{End of Quiz}
\end{center}

\end{document}